
\documentclass{article}
\usepackage{colortbl}
\usepackage{makecell}
\usepackage{multirow}
\usepackage{supertabular}

\begin{document}

\newcounter{utterance}

\centering \large Interaction Transcript for game `cladder', experiment `full\_v1.5\_default', episode 2646 with refineonly{-}Qwen3.5{-}4B.
\vspace{24pt}

{ \footnotesize  \setcounter{utterance}{1}
\setlength{\tabcolsep}{0pt}
\begin{supertabular}{c@{$\;$}|p{.15\linewidth}@{}p{.15\linewidth}p{.15\linewidth}p{.15\linewidth}p{.15\linewidth}p{.15\linewidth}}
    \# & \multicolumn{2}{c}{Player} && \multicolumn{2}{c}{Game Master} \\
    \hline

    \theutterance \stepcounter{utterance}  
    & & & \multicolumn{4}{p{0.6\linewidth}}{
        \cellcolor[rgb]{0.9,0.9,0.9}{
            \makecell[{{p{\linewidth}}}]{
                \texttt{\tiny{[P1$\langle$GM]}}
                \texttt{You are an expert in causal inference. The following question is not a typical commonsense query, but rather a meticulously designed question created by a professor specializing in causal inference, intended to assess the students' mastery of the course content.} \\
\texttt{Imagine a self{-}contained, hypothetical world with only the following conditions, and without any unmentioned factors or causal relationships: Gender has a direct effect on smoking and lung cancer. Smoking has a direct effect on tar deposit. Tar deposit has a direct effect on lung cancer. Gender is unobserved. For nonsmokers and with no tar deposit, the probability of lung cancer is 79\%. For nonsmokers and with high tar deposit, the probability of lung cancer is 49\%. For smokers and with no tar deposit, the probability of lung cancer is 71\%. For smokers and with high tar deposit, the probability of lung cancer is 40\%. For nonsmokers, the probability of high tar deposit is 79\%. For smokers, the probability of high tar deposit is 40\%. For smokers, would it be more likely to see lung cancer if the person had been a nonsmoker?} \\
\texttt{Start your answer with 'Yes' or 'No', followed by additional reasoning or evidence to support your explanation.} \\
            }
        }
    }
     \\ \\

    \theutterance \stepcounter{utterance}  
    & \multicolumn{4}{p{0.6\linewidth}}{
        \cellcolor[rgb]{0.561,0.737,0.561}{
            \makecell[{{p{\linewidth}}}]{
                \texttt{\tiny{[P1$\rangle$GM]}}
                \texttt{Yes} \\
\\ 
\texttt{**Reasoning:**} \\
\\ 
\texttt{To determine whether a smoker would be more likely to see lung cancer if they had been a nonsmoker, we must compare the counterfactual probability of lung cancer for a smoker under the intervention \$do(\textbackslash{}text\{Nonsmoker\})\$ against the observed probability for a smoker, or more simply, compare the marginal probability of lung cancer for the group of nonsmokers versus the group of smokers in this hypothetical world. The question asks if the likelihood decreases when switching from the "Smoker" state to the "Nonsmoker" state.} \\
\\ 
\texttt{1.  **Calculate the Marginal Probability of Lung Cancer for Nonsmokers (\$P(LC | \textbackslash{}text\{Nonsmoker\})\$):**} \\
\texttt{    *   We are given the conditional probabilities of lung cancer based on tar deposit status for nonsmokers:} \\
\texttt{        *   \$P(LC | \textbackslash{}text\{Nonsmoker, High Tar\}) = 49\textbackslash{}\%\$} \\
\texttt{        *   \$P(LC | \textbackslash{}text\{Nonsmoker, No Tar\}) = 79\textbackslash{}\%\$} \\
\texttt{    *   We are given the probability of high tar deposit for nonsmokers:} \\
\texttt{        *   \$P(\textbackslash{}text\{High Tar\} | \textbackslash{}text\{Nonsmoker\}) = 79\textbackslash{}\%\$} \\
\texttt{        *   \$P(\textbackslash{}text\{No Tar\} | \textbackslash{}text\{Nonsmoker\}) = 100\textbackslash{}\% {-} 79\textbackslash{}\% = 21\textbackslash{}\%\$} \\
\texttt{    *   Using the law of total probability:} \\
\texttt{        \$\$P(LC | \textbackslash{}text\{Nonsmoker\}) = P(LC | \textbackslash{}text\{Nonsmoker, High Tar\}) \textbackslash{}times P(\textbackslash{}text\{High Tar\} | \textbackslash{}text\{Nonsmoker\}) + P(LC | \textbackslash{}text\{Nonsmoker, No Tar\}) \textbackslash{}times P(\textbackslash{}text\{No Tar\} | \textbackslash{}text\{Nonsmoker\})\$\$} \\
\texttt{        \$\$P(LC | \textbackslash{}text\{Nonsmoker\}) = 0.49 \textbackslash{}times 0.79 + 0.79 \textbackslash{}times 0.21\$\$} \\
\texttt{        \$\$P(LC | \textbackslash{}text\{Nonsmoker\}) = 0.3871 + 0.1659 = 0.553 \textbackslash{}text\{ (or \} 55.3\textbackslash{}\%)\$\$} \\
\\ 
\texttt{2.  **Calculate the Marginal Probability of Lung Cancer for Smokers (\$P(LC | \textbackslash{}text\{Smoker\})\$):**} \\
\texttt{    *   We are given the conditional probabilities of lung cancer based on tar deposit status for smokers:} \\
\texttt{        *   \$P(LC | \textbackslash{}text\{Smoker, High Tar\}) = 40\textbackslash{}\%\$} \\
\texttt{        *   \$P(LC | \textbackslash{}text\{Smoker, No Tar\}) =} \\
            }
        }
    }
    & & \\ \\

    \theutterance \stepcounter{utterance}  
    & & & \multicolumn{2}{p{0.3\linewidth}}{
        \cellcolor[rgb]{0.9,0.9,0.9}{
            \makecell[{{p{\linewidth}}}]{
                \texttt{\tiny{[GM$|$GM]}}
                \texttt{Yes} \\
            }
        }
    }
    & & \\ \\

    \theutterance \stepcounter{utterance}  
    & & & \multicolumn{2}{p{0.3\linewidth}}{
        \cellcolor[rgb]{0.9,0.9,0.9}{
            \makecell[{{p{\linewidth}}}]{
                \texttt{\tiny{[GM$|$GM]}}
                \texttt{no} \\
            }
        }
    }
    & & \\ \\

    \theutterance \stepcounter{utterance}  
    & & & \multicolumn{2}{p{0.3\linewidth}}{
        \cellcolor[rgb]{0.9,0.9,0.9}{
            \makecell[{{p{\linewidth}}}]{
                \texttt{\tiny{[GM$|$GM]}}
                \texttt{game\_result = LOSE} \\
            }
        }
    }
    & & \\ \\

\end{supertabular}
}

\end{document}
